\documentclass[12pt]{article}
\usepackage[utf8]{inputenc}
\usepackage[T1]{fontenc}
\usepackage[margin=1in]{geometry}
\usepackage{setspace}
\usepackage{xcolor}
\usepackage[colorlinks=true, urlcolor=blue!60!black, linkcolor=blue!60!black]{hyperref}
\setlength{\parindent}{0em}
\setlength{\parskip}{0.6em}
\renewcommand{\baselinestretch}{1.0}
\begin{document}
\begin{center}
{\LARGE \textit{United States v. Virginia}, 518 U.S. 515 (1996)\par}
\vspace{0.5em}
{\large Case Brief --- Two-Page Edition\par}
\end{center}

\section{Memory Jogger}
VMI's exclusion of women lacked an ``exceedingly persuasive justification,'' and a separate, inferior women's program did not cure it (\href{https://supreme.justia.com/cases/federal/us/518/515/}{518 U.S. 515}).

\section{Detailed Case Facts}
VMI, Virginia's only single-sex public college, trained male ``citizen-soldiers'' through a grueling ``adversative method'' and had a vast endowment and alumni network; it ignored inquiries from 347 women. After the United States sued and lost on appeal, Virginia created VWIL at private Mary Baldwin College, using a ``cooperative method'' with fewer resources, weaker faculty, and no engineering degree.

\section{Procedural History}
The district court upheld VMI's policy; the Fourth Circuit vacated; on remand the lower courts approved VWIL under a ``substantive comparability'' test.

\section{Judicial Votes}
7--1. Ginsburg wrote for Stevens, O'Connor, Kennedy, Souter, and Breyer; Rehnquist concurred in the judgment; Scalia dissented; Thomas did not participate.

\section{Holding}
A State must show an exceedingly persuasive justification for a sex classification: important objectives, substantially related means, a genuine rather than \textit{post hoc} rationale, and no overbroad generalizations. VWIL, unequal in resources and prestige, was no remedy.

\section{Analysis of Opinions}
\textbf{Majority:} VMI was not founded for diversity, and fears about women were ``self-fulfilling prophec[ies]''; some women could meet every demand. Under \textit{Sweatt v.\ Painter}, VWIL was a ``pale shadow.'' \textbf{Rehnquist:} the test is \textit{Craig v.\ Boren}'s; a truly equal women's school would have sufficed. \textbf{Scalia:} tradition should be changed democratically; the Court covertly adopted strict scrutiny and doomed public single-sex education.

\section{Examples: Future Applications}
\textbf{Same-Side:} a male-only state maritime academy citing women's average preferences fails. \textbf{Opposite-Side:} an engineering fellowship for underrepresented women, backed by data, can survive as compensatory. \textbf{Fence-Sitter:} twin, equally funded boys' and girls' public high schools, the question footnote 7 reserved.

\section{Critique}
Conservatives call the standard ``strict scrutiny in disguise,'' with Scalia's dissent the ``locus classicus'' of that view. Cass Sunstein praised the opinion as ``ambitious'' and anti-hierarchical, while Katharine Bartlett faulted it for accepting VMI's hyper-masculine method as the baseline.

\section{Key Quotations}
\begin{quote}
``The justification must be genuine, not hypothesized or invented post hoc in response to litigation. And it must not rely on overbroad generalizations about the different talents, capacities, or preferences of males and females.''
\end{quote}
\noindent --- Justice Ginsburg, Opinion of the Court, 518 U.S. at 533.

\end{document}
